\textbf{What the evidence establishes.} Source diversity, executed
mechanisms, and one-run correctness differences are distinct measurements.
The Phase-I pattern did not recur in the revised Phase-II setting; multiple
protocol changes prevent identifying the cause. The corrected analysis of
the original gate shows no clear headroom improvement. Because that gate
also excludes allowed prompt strategies, it does not identify the effect
of verification alone.

\textbf{What remains unidentified.} Same-code reruns can disagree and
create oracle gains. Outcome flip rates measure instability, not the
standard error of a treatment contrast, and cannot define a numerical
noise floor for that contrast. A matched-cost clone control, independent
repeat validation, and calibration covering all admitted mechanism types
are needed before attributing headroom to stable additional capability.

\textbf{Limitations.} The main comparison uses one text-to-SQL target and
six fixed builders, with three generation seeds per builder. Phase I is
discovery evidence with disclosed task overlap, a proxy judge, and a
post-trace human control. The new statistical analysis is post-review and
does not add independent executions. K-matched and R2 results now share the
corrected data and resampling chain. W1/W3 now share the canonical data
but remain descriptive; selection uncertainty and independent utility
validation are outstanding.
The completed one-seed MATH-500 free-arm probe
(Appendix~\ref{app:crossdomain}) extends the observed outcome comparison
to a second domain, but does not supply a clone control, a repaired-gate
factorial, or an independent test of stable complementarity.
